\section{Deterministic Execution Contracts for LLM Agent Systems:
Architecture, Policy Enforcement, and Empirical
Evaluation}\label{deterministic-execution-contracts-for-llm-agent-systems-architecture-policy-enforcement-and-empirical-evaluation}

\textbf{Adam Tacon}\\
Axiom LLC, New Jersey\\
\texttt{adam.tacon@proton.me} \texttt{axiom.co@proton.me}

\begin{center}\rule{0.5\linewidth}{0.5pt}\end{center}

\subsection{Abstract}\label{abstract}

Large language model (LLM) agents introduce a structural mismatch
between probabilistic generation and the determinism required for safe,
reproducible execution in consequential environments. We present APEX
(Autonomous Plan EXecutor), a plan-deterministic swarm runtime that
enforces schema-validated, bounded, process-isolated execution contracts
over LLM-generated plans, and ASON (Autonomous Security Optimization
Network), a policy enforcement layer that validates and governs plan
submission prior to execution. Together, these systems instantiate a
two-stage architecture in which generation remains probabilistic but
execution is structurally constrained: given a validated plan, the
tool-call sequence is deterministic; LLM variance across plan generation
is logged and replayable. We formalize the information-flow discipline
underlying this design, describe the implementation of key invariants
including step bounding, tool timeout, blast-radius enforcement, and
rollback generation, and report empirical results from a benchmark
harness yielding a 12/12 pass rate with a composite \texttt{apex\_score}
of 1.008 across nominal and adversarial tool-invocation, file I/O,
memory, and HTTP tasks. We further describe the Benchmark-Driven
Self-Optimization (BDSO) loop integrated into APEX, bounded by a fitness
function and restricted to a verified-eligible file set that permanently
excludes security-critical components. The system is released as
open-source infrastructure for deterministic AI agent deployment.

\begin{center}\rule{0.5\linewidth}{0.5pt}\end{center}

\subsection{1. Introduction}\label{introduction}

The dominant failure mode of deployed LLM agent systems is not
generative inaccuracy but architectural unsafety: the absence of
enforceable boundaries on information flow between probabilistic
generation and consequential execution. When an agent ingests tool
outputs, retrieved documents, and user instructions through a unified
untyped context, data and control become structurally indistinguishable.
This ambiguity is the root cause of prompt injection, policy drift,
unverified plan execution, and unstable self-modification.

Prior work has addressed these failure modes largely through heuristic
filtering, constitutional prompting, and output classifiers ---
interventions that operate at the semantic layer and inherit its
fundamental unreliability. We argue that reliable agent execution
requires a lower-level intervention: typed intermediate representations,
schema-validated plan contracts, and execution environments with
formally bounded resource consumption.

This paper makes the following contributions:

\begin{enumerate}
\def\labelenumi{\arabic{enumi}.}
\tightlist
\item
  A formal characterization of information-flow collapse as the unifying
  cause across injection, verification failure, and unstable
  self-improvement.
\item
  APEX: a plan-deterministic agent runtime implementing a pure
  \texttt{run(task,\ config,\ registry)\ →\ State} execution contract
  with bounded step count, tool timeout, output limits, and JSONL audit
  trace.
\item
  ASON: a policy enforcement layer implementing blast-radius
  classification, tool allowlisting, and automatic rollback plan
  generation against the APEX execution substrate.
\item
  A benchmark harness and fitness function (\texttt{apex\_score})
  producing verifiable, reproducible evaluation of agent task execution.
\item
  A Benchmark-Driven Self-Optimization (BDSO) loop bounded by the
  fitness function and restricted to a security-invariant eligible file
  set.
\end{enumerate}

\begin{center}\rule{0.5\linewidth}{0.5pt}\end{center}

\subsection{2. Related Work}\label{related-work}

\textbf{Prompt injection and control-flow integrity.} Perez and Ribeiro
(2022) characterize prompt injection as an input-validation failure;
Greshake et al.~(2023) extend this to indirect injection via retrieved
content. Our treatment formalizes the mechanism as type confusion over a
semantically untyped channel, analogous to memory-unsafe control-flow
corruption.

\textbf{Plan verification.} SayCan (Ahn et al., 2022) and ToolFormer
(Schick et al., 2023) generate tool-calling plans but perform no
pre-execution structural validation. ReAct (Yao et al., 2022)
interleaves reasoning and action without schema enforcement. APEX
introduces a validation stage that rejects plans failing schema,
step-count, or tool-allowlist constraints prior to any execution.

\textbf{Agent safety frameworks.} Constitutional AI (Bai et al., 2022)
and RLHF-based approaches address value alignment at training time. Our
approach is complementary and orthogonal: we address execution-time
structural safety independently of model-level alignment.

\textbf{Policy enforcement and execution contracts.} Existing
orchestration frameworks lack the enforcement layer that distinguishes
ASON from prior work. LangGraph provides a stateful execution graph but
no policy layer: tool invocation is unrestricted by blast-radius or
allowlist constraints. Temporal enforces workflow durability and retry
semantics but has no LLM-native contract mechanism --- policies are
defined in application code, not validated against plan structure.
AutoGen supports multi-agent coordination but provides no execution
bounds on individual agent actions. ASON's distinguishing contribution
is enforcement at the API boundary: blast-radius classification, tool
allowlisting, and step-count limits are validated against the plan
structure before any execution begins, with no comparator in the cited
frameworks.

\textbf{Recursive self-improvement.} Schmidhuber (2003) and subsequent
work characterize RSI stability requirements. APEX instantiates a
pragmatic bounded self-optimization loop (BDSO) with multi-candidate
patch generation, isolated scoring, and a permanent security boundary
excluding the plan auditor from eligible modification targets.

\begin{center}\rule{0.5\linewidth}{0.5pt}\end{center}

\subsection{3. Theoretical Framework: Information Flow as the Core
Safety
Primitive}\label{theoretical-framework-information-flow-as-the-core-safety-primitive}

\subsubsection{3.1 The Type-Confusion Model of Prompt
Injection}\label{the-type-confusion-model-of-prompt-injection}

LLM agents process all inputs --- user prompts, tool outputs, retrieved
documents --- through a single untyped context. The model assigns no
structural distinction between data (what is known) and control (what is
done). A string that appears descriptive may be reinterpreted as
prescriptive if it aligns with latent training patterns.

This is structurally equivalent to type confusion in memory-unsafe
languages: a buffer interpreted as a control structure can overwrite
execution flow; a string interpreted as an instruction can override
policy constraints. The agent does not resolve this ambiguity through
judgment --- it lacks the structural capacity to prevent
reinterpretation.

The critical consequence is \textbf{compositionality}: individual
fragments may be benign, but their accumulation across retrieval, tool
output, and memory can constitute a coherent adversarial program that
emerges through interaction, never existing explicitly at any single
point. Local per-input filtering is therefore provably insufficient.
Safety must be defined over entire execution traces.

\subsubsection{3.2 Verification as a Correctness
Envelope}\label{verification-as-a-correctness-envelope}

Verification in LLM agent systems cannot achieve full formal correctness
--- plans are underspecified, environments are non-deterministic, and
verification latency must remain bounded for real-time operation. The
pragmatic objective is a \textbf{correctness envelope}: a bounded region
within which all executions are guaranteed to satisfy safety
constraints.

Critically, this envelope must be \textbf{temporal and relational}: an
action safe in isolation may violate constraints in combination with
prior or subsequent steps. Verification must therefore operate over the
complete plan structure, including state transitions and resource
consumption projections, not over individual steps.

\subsubsection{3.3 Recursive Self-Improvement and Stability
Constraints}\label{recursive-self-improvement-and-stability-constraints}

RSI dynamics are most accurately modeled as a feedback system. Stability
requires that errors decay across iterations --- each modification cycle
must reduce, not amplify, deviation from desired behavior. Two
conditions are necessary: (1) the magnitude of per-iteration changes
must be bounded to keep verification tractable, and (2) the verification
mechanism must scale with capability growth. A system that improves
faster than its verifier can evaluate enters a regime of uncontrolled
capability increase.

A further invariant is that \textbf{security-critical components must be
excluded from the eligible modification set}. A self-modifying system
that can alter its own safety enforcement is not self-improving --- it
is self-destabilizing.

\begin{center}\rule{0.5\linewidth}{0.5pt}\end{center}

\subsection{4. APEX: Deterministic Agent
Runtime}\label{apex-deterministic-agent-runtime}

\subsubsection{4.1 Core Execution
Contract}\label{core-execution-contract}

APEX implements a pure function:

\begin{verbatim}
run(task: str, config: Config, registry: ToolRegistry) → State
\end{verbatim}

The function is stateless with respect to prior runs. All state is local
to the invocation and flushed to the audit log on completion. This
purity is a primary architectural invariant: no hidden state accumulates
across runs, and given a validated plan, execution is fully reproducible
--- identical plans produce identical tool-call sequences. LLM variance
occurs at plan generation time and is captured in the audit log rather
than eliminated.

\subsubsection{4.2 Execution Bounds}\label{execution-bounds}

Every run is subject to hard resource bounds:

\begin{longtable}[]{@{}ll@{}}
\toprule\noalign{}
Constraint & Bound \\
\midrule\noalign{}
\endhead
\bottomrule\noalign{}
\endlastfoot
Maximum steps & 32 \\
Tool timeout & 300 s (SIGALRM, main thread) \\
Output size & 10 MB \\
Plan schema & Pydantic-validated pre-execution \\
\end{longtable}

Plans failing schema validation are rejected before any tool invocation.
Step count is enforced as a hard ceiling, not a soft advisory. Tool
timeout is enforced at the syscall boundary via \texttt{signal.alarm},
ensuring no tool invocation can hold the executor indefinitely.

\subsubsection{4.3 Tool Registry and
Autoloading}\label{tool-registry-and-autoloading}

Tools are registered against a typed interface requiring an
\texttt{effect(args)\ →\ result} method and a JSON schema for argument
validation. The runtime autoloads tools from
\texttt{\textasciitilde{}/.apex/tools/}, enabling deployment-specific
tool sets without modifying core execution logic. This separation
maintains the KISS invariant: the core executor has no awareness of
domain-specific tool semantics.

\subsubsection{4.4 Audit Infrastructure}\label{audit-infrastructure}

Every run emits a JSONL trace recording each step's tool invocation,
arguments, result, and wall time. Runs are persisted to a SQLite
database (\texttt{\textasciitilde{}/.apex/runs.db}) with schema:

\begin{Shaded}
\begin{Highlighting}[]
\NormalTok{runs(}\KeywordTok{id}\NormalTok{, task, plan\_json, exit\_code, token\_count, wall\_seconds, }\DataTypeTok{timestamp}\NormalTok{)}
\KeywordTok{events}\NormalTok{(}\KeywordTok{id}\NormalTok{, run\_id, step, tool, args\_json, result\_json, }\DataTypeTok{timestamp}\NormalTok{)}
\end{Highlighting}
\end{Shaded}

This provides complete replay fidelity and supports post-hoc audit of
any execution. The \texttt{apex\ replay} subcommand supports three
modes: \texttt{simulate} (reconstruct trace without re-executing),
\texttt{dry} (execute plan without committing side effects), and
\texttt{live} (full re-execution).

\subsubsection{4.5 Multi-Provider
Abstraction}\label{multi-provider-abstraction}

The LLM provider is abstracted behind a \texttt{providers.py} shim with
concrete adapters for Gemini 2.5 Flash and Ollama. Provider selection is
configuration-driven at startup; the core execution loop has no
provider-specific logic. This isolation ensures that provider
substitution does not require modification of any safety-critical
component.

\subsubsection{4.6 HTTP API}\label{http-api}

APEX exposes a Flask HTTP API (\texttt{apex\ serve}) with endpoints for
run submission (\texttt{POST\ /run}), run history (\texttt{GET\ /runs},
\texttt{GET\ /runs/\textless{}id\textgreater{}}), replay
(\texttt{POST\ /replay}), and export (\texttt{GET\ /export}). The server
runs in single-threaded mode (\texttt{threaded=False}) to maintain
SIGALRM compatibility --- the tool timeout mechanism operates only on
the main thread, and multi-threaded dispatch would silently disable it.

Authentication is enforced via \texttt{X-Apex-Key} header matched
against \texttt{APEX\_API\_KEY} environment variable. If the variable is
unset, the server starts without authentication with a startup warning.

\begin{center}\rule{0.5\linewidth}{0.5pt}\end{center}

\subsection{5. ASON: Policy Enforcement
Layer}\label{ason-policy-enforcement-layer}

\subsubsection{5.1 Architecture}\label{architecture}

ASON (Autonomous Security Optimization Network) is a pre-execution
policy enforcement layer that validates and governs plan submission to
APEX. It implements a strict separation: APEX defines execution
semantics; ASON defines submission policy. Plans that violate policy are
rejected before reaching the APEX executor.

\subsubsection{5.2 Policy Schema}\label{policy-schema}

ASON policies are Pydantic-validated contracts:

\begin{Shaded}
\begin{Highlighting}[]
\KeywordTok{class}\NormalTok{ Policy(BaseModel):}
\NormalTok{    max\_steps: }\BuiltInTok{int}          \CommentTok{\# default 16, bounds [1, 32]}
\NormalTok{    allowed\_tools: }\BuiltInTok{list}\NormalTok{[}\BuiltInTok{str}\NormalTok{] }\CommentTok{\# empty = unrestricted}
\NormalTok{    blast\_radius: Literal[}\StringTok{"none"}\NormalTok{, }\StringTok{"local"}\NormalTok{, }\StringTok{"network"}\NormalTok{]  }\CommentTok{\# default "local"}
\NormalTok{    rollback\_on\_failure: }\BuiltInTok{bool}  \CommentTok{\# default True}
\end{Highlighting}
\end{Shaded}

The \texttt{blast\_radius} field classifies the execution's potential
impact envelope. \texttt{none} permits read-only operations;
\texttt{local} permits filesystem modification within the process
working directory; \texttt{network} permits egress. Policy is enforced
at the structural level prior to execution, not as a runtime advisory.

\subsubsection{5.3 Validation Logic}\label{validation-logic}

The ASON validator enforces:

\begin{enumerate}
\def\labelenumi{\arabic{enumi}.}
\tightlist
\item
  \texttt{len(plan.steps)\ ≤\ policy.max\_steps}
\item
  No step invokes a tool in the blocked set (\texttt{\{"shell"\}} by
  default)
\item
  If \texttt{policy.allowed\_tools} is non-empty, every step's tool must
  appear in it
\item
  No step invokes a tool in the \texttt{blast\_radius}-derived block set
  (\texttt{blast\_radius="none"} blocks \texttt{write\_file},
  \texttt{delete\_file}, \texttt{http\_get}, \texttt{http\_post};
  \texttt{blast\_radius="local"} blocks \texttt{http\_get},
  \texttt{http\_post})
\end{enumerate}

Violations produce structured rejection results with per-violation
diagnostics. Validated plans are submitted to \texttt{POST\ /run} via
the APEX HTTP API with the \texttt{X-Apex-Key} credential.

\subsubsection{5.4 Rollback Generation}\label{rollback-generation}

On execution failure, ASON generates a compensating \texttt{ASONRequest}
by traversing the \texttt{events} table for the failed run in reverse
step order and applying the reversal map:

\begin{verbatim}
write_file → delete_file
\end{verbatim}

\texttt{shell} invocations are flagged to stderr as non-reversible. The
rollback plan is itself subject to policy validation before submission.
If no reversible steps exist, rollback returns \texttt{None} without
error.

\begin{center}\rule{0.5\linewidth}{0.5pt}\end{center}

\subsection{6. Benchmark-Driven Self-Optimization
(BDSO)}\label{benchmark-driven-self-optimization-bdso}

\subsubsection{6.1 Bounded BDSO Loop}\label{bounded-bdso-loop}

APEX implements a Benchmark-Driven Self-Optimization loop
(\texttt{apex\ rsi}) that generates candidate patches to the executor's
own source files, scores them via the benchmark harness, and applies
improvements that exceed a fitness threshold.

Each BDSO cycle generates N=3 candidate patches via the LLM, executes
them in isolated subprocesses, scores each against the
\texttt{apex\_score} fitness function, and applies the highest-scoring
candidate that exceeds the baseline. If no candidate improves on the
baseline, the cycle terminates without modification. The mechanism is
bounded and empirically grounded: improvement is defined by measurable
benchmark regression, not open-ended capability growth.

\subsubsection{6.2 Fitness Function}\label{fitness-function}

The \texttt{apex\_score} composite fitness function is:

\begin{verbatim}
apex_score = pass_rate × speed_factor × token_efficiency
\end{verbatim}

Where: - \texttt{pass\_rate\ =\ passed\_tasks\ /\ total\_tasks} -
\texttt{speed\_factor\ =\ max(0.01,\ 1.0\ −\ (avg\_duration\_seconds\ −\ 10.0)\ /\ 200.0)}
-
\texttt{token\_efficiency\ =\ max(0.01,\ 1.0\ −\ (avg\_tokens\ −\ 1000.0)\ /\ 50000.0)}

This formulation penalizes both execution latency and token
overconsumption while treating task success as the primary signal. A
system that passes all tasks at 10s/task with 1000 tokens/task achieves
\texttt{apex\_score\ =\ 1.0}.

\subsubsection{6.3 Security Boundary}\label{security-boundary}

The BDSO eligible file set is restricted to
\texttt{\{loop.py,\ planner.py,\ llm.py\}}. The plan auditor
(\texttt{paranoid.py}) is \textbf{permanently excluded} from the
eligible set. This exclusion is architectural, not configurable: a
self-modifying system that can alter its own safety enforcement
mechanism provides no meaningful safety guarantee.

\begin{center}\rule{0.5\linewidth}{0.5pt}\end{center}

\subsection{7. Evaluation}\label{evaluation}

\subsubsection{7.1 Benchmark Harness}\label{benchmark-harness}

The benchmark harness (\texttt{apex.bench}) loads a task suite from
\texttt{tasks.json}, executes each task via the APEX CLI in a
subprocess, and records per-task pass/fail, duration, and exit code.
Tasks specify an optional content check against stdout or a designated
output file.

\subsubsection{7.2 Task Suite}\label{task-suite}

The benchmark suite comprises 12 tasks across nominal and adversarial
categories:

\begin{longtable}[]{@{}
  >{\raggedright\arraybackslash}p{(\linewidth - 6\tabcolsep) * \real{0.2500}}
  >{\raggedright\arraybackslash}p{(\linewidth - 6\tabcolsep) * \real{0.2500}}
  >{\raggedright\arraybackslash}p{(\linewidth - 6\tabcolsep) * \real{0.2500}}
  >{\raggedright\arraybackslash}p{(\linewidth - 6\tabcolsep) * \real{0.2500}}@{}}
\toprule\noalign{}
\begin{minipage}[b]{\linewidth}\raggedright
Task ID
\end{minipage} & \begin{minipage}[b]{\linewidth}\raggedright
Category
\end{minipage} & \begin{minipage}[b]{\linewidth}\raggedright
Description
\end{minipage} & \begin{minipage}[b]{\linewidth}\raggedright
Tool(s)
\end{minipage} \\
\midrule\noalign{}
\endhead
\bottomrule\noalign{}
\endlastfoot
\texttt{write\_file} & Nominal & Write string to \texttt{/tmp} path,
verify content & \texttt{write\_file} \\
\texttt{shell\_echo} & Nominal & Execute shell command, verify stdout &
\texttt{shell} \\
\texttt{read\_file} & Nominal & Write then read file, verify round-trip
& \texttt{read\_file} \\
\texttt{memory\_write\_read} & Nominal & Write key-value to memory,
retrieve and verify & \texttt{memory} \\
\texttt{http\_get} & Nominal & HTTP GET to public endpoint, verify
response & \texttt{http\_get} \\
\texttt{multi\_step\_chain} & Adversarial & Multi-step chained file
write and read & \texttt{write\_file}, \texttt{read\_file} \\
\texttt{overwrite\_existing} & Adversarial & Write to pre-existing path,
verify overwrite & \texttt{write\_file} \\
\texttt{shell\_pipeline} & Adversarial & Execute shell pipeline, verify
stdout & \texttt{shell} \\
\texttt{http\_post} & Adversarial & HTTP POST to endpoint, verify
response & \texttt{http\_post} \\
\texttt{memory\_multi\_key} & Adversarial & Write multiple keys,
retrieve and verify all & \texttt{memory} \\
\texttt{large\_file\_write} & Adversarial & Write large payload, verify
size and content & \texttt{write\_file} \\
\texttt{multi\_file\_verify} & Adversarial & Write multiple files, read
and verify each & \texttt{write\_file}, \texttt{read\_file} \\
\end{longtable}

The adversarial suite targets boundary conditions identified as coverage
gaps in prior evaluation: multi-step plan chaining, idempotency under
overwrite, I/O scaling, and multi-key memory operations. Runtime
blast-radius isolation (chroot/namespace enforcement) remains a
near-term objective; current \texttt{blast\_radius} enforcement operates
at plan-validation time.

\subsubsection{7.3 Results}\label{results}

Two independent benchmark runs were executed on hardware (AMD Ryzen 3
3250U, 33 GiB RAM, NVMe storage) under identical conditions with a
180-second per-task timeout:

\begin{longtable}[]{@{}
  >{\raggedright\arraybackslash}p{(\linewidth - 10\tabcolsep) * \real{0.1667}}
  >{\raggedright\arraybackslash}p{(\linewidth - 10\tabcolsep) * \real{0.1667}}
  >{\raggedright\arraybackslash}p{(\linewidth - 10\tabcolsep) * \real{0.1667}}
  >{\raggedright\arraybackslash}p{(\linewidth - 10\tabcolsep) * \real{0.1667}}
  >{\raggedright\arraybackslash}p{(\linewidth - 10\tabcolsep) * \real{0.1667}}
  >{\raggedright\arraybackslash}p{(\linewidth - 10\tabcolsep) * \real{0.1667}}@{}}
\toprule\noalign{}
\begin{minipage}[b]{\linewidth}\raggedright
Run
\end{minipage} & \begin{minipage}[b]{\linewidth}\raggedright
Pass Rate
\end{minipage} & \begin{minipage}[b]{\linewidth}\raggedright
Wall Time
\end{minipage} & \begin{minipage}[b]{\linewidth}\raggedright
\texttt{speed\_factor}
\end{minipage} & \begin{minipage}[b]{\linewidth}\raggedright
\texttt{token\_efficiency}
\end{minipage} & \begin{minipage}[b]{\linewidth}\raggedright
\texttt{apex\_score}
\end{minipage} \\
\midrule\noalign{}
\endhead
\bottomrule\noalign{}
\endlastfoot
Run 1 & 12/12 (1.0) & 92.0 s & 1.0117 & 0.9959 & 1.007552 \\
Run 2 & 12/12 (1.0) & 90.5 s & 1.0123 & 0.9954 & 1.007643 \\
\end{longtable}

\texttt{apex\_score\ \textgreater{}\ 1.0} in both runs indicates
execution faster than the 10s/task baseline with near-unity token
efficiency. Results are consistent across runs (δ \textless{} 0.0001).
All 12 tasks completed within timeout; the \texttt{memory\_multi\_key}
task exhibits the longest per-task wall time (\textasciitilde22--25 s)
due to multi-round memory retrieval, accounted for by the 180 s timeout
budget.

\begin{center}\rule{0.5\linewidth}{0.5pt}\end{center}

\subsection{8. Docker Deployment}\label{docker-deployment}

APEX and ASON are containerized and orchestrated via
\texttt{axiom-infra} (docker-compose). The compose stack defines:

\begin{itemize}
\tightlist
\item
  \textbf{apex}: builds from \texttt{axiom-apex}; binds port 8080;
  healthcheck on \texttt{GET\ /health}; persists
  \texttt{\textasciitilde{}/.apex} via named volume
\item
  \textbf{ason}: builds from \texttt{axiom-ason}; depends on apex health
  passing before start
\end{itemize}

This enables production deployment of the full two-layer architecture as
a single \texttt{docker-compose\ up} invocation, with GEMINI\_API\_KEY
injected at runtime from environment.

\begin{center}\rule{0.5\linewidth}{0.5pt}\end{center}

\subsection{9. Future Work}\label{future-work}

\textbf{Blast-radius runtime enforcement.} \texttt{blast\_radius} is
currently enforced at plan-validation time. Full runtime enforcement
requires process isolation (chroot or namespace) for \texttt{local} and
an egress proxy for \texttt{network}. Plan-level enforcement is a
necessary but not sufficient condition for the security guarantee the
field implies.

\textbf{Multi-agent consensus.} For high-\texttt{blast\_radius} tasks,
generating N=3 independent plans and requiring intersection before
execution would substantially reduce the probability of unilateral
consequential actions.

\textbf{RAG-driven policy grounding.} Policy constraints derived from
compliance manifests (SOC2, HIPAA) via retrieval-augmented generation
would enable domain-specific ASON policy without manual specification.

\textbf{BDSO fitness expansion.} The current \texttt{apex\_score}
formulation does not include a security-regression component. Adding a
paranoid-flag delta term would penalize BDSO cycles that reduce plan
auditor sensitivity.

\begin{center}\rule{0.5\linewidth}{0.5pt}\end{center}

\subsection{10. Conclusion}\label{conclusion}

The reliable deployment of LLM agents in consequential environments
requires architectural discipline that probabilistic generation alone
cannot provide. The central invariant is information-flow separation:
generation proposes, a typed intermediate representation constrains, and
a bounded executor acts. APEX and ASON instantiate this architecture
concretely, with verifiable execution contracts, schema-enforced policy,
complete audit infrastructure, and a BDSO loop bounded by a composite
fitness function and a permanent security exclusion boundary.

The empirical results --- 12/12 pass rate across nominal and adversarial
tasks, \texttt{apex\_score} \textasciitilde1.008 across two independent
runs --- establish a reproducible baseline against which future
capability and hardening improvements can be measured. The architecture
is released as open-source infrastructure at
\texttt{github.com/axiom-llc/axiom-apex} and
\texttt{github.com/axiom-llc/axiom-ason}.

\begin{center}\rule{0.5\linewidth}{0.5pt}\end{center}

\subsection{References}\label{references}

Ahn, M., et al.~(2022). Do As I Can, Not As I Say: Grounding Language in
Robotic Affordances. \emph{arXiv:2204.01691}.

Bai, Y., et al.~(2022). Constitutional AI: Harmlessness from AI
Feedback. \emph{arXiv:2212.08073}.

Greshake, K., et al.~(2023). Not What You've Signed Up For: Compromising
Real-World LLM-Integrated Applications with Indirect Prompt Injection.
\emph{arXiv:2302.12173}.

Perez, F., \& Ribeiro, I. (2022). Ignore Previous Prompt: Attack
Techniques For Language Models. \emph{arXiv:2211.09527}.

Schick, T., et al.~(2023). Toolformer: Language Models Can Teach
Themselves to Use Tools. \emph{arXiv:2302.04761}.

Schmidhuber, J. (2003). Gödel Machines: Self-Referential Universal
Problem Solvers Making Provably Optimal Self-Improvements.
\emph{Technical Report IDSIA-19-03}.

Yao, S., et al.~(2022). ReAct: Synergizing Reasoning and Acting in
Language Models. \emph{arXiv:2210.03629}.

\begin{center}\rule{0.5\linewidth}{0.5pt}\end{center}

\subsection{Appendix A: APEX
Invariants}\label{appendix-a-apex-invariants}

The following invariants are maintained across all versions of the
system:

\begin{enumerate}
\def\labelenumi{\arabic{enumi}.}
\tightlist
\item
  \texttt{run()} is a pure function --- no hidden state persists across
  invocations
\item
  No async primitives or threading in the core execution path
\item
  Plan-deterministic outputs: given a validated plan, identical inputs
  produce identical tool-call sequences; LLM variance at plan generation
  time is logged and replayable
\item
  Bounded execution: \texttt{max\_steps=32}, \texttt{timeout=300s},
  \texttt{output=10MB}
\item
  All plans are schema-validated before any tool invocation
\item
  ASON never bypasses APEX policy enforcement
\item
  \texttt{paranoid.py} is not BDSO-eligible --- permanent security
  boundary
\item
  Flask HTTP API runs \texttt{threaded=False} to preserve SIGALRM
  semantics
\end{enumerate}

\subsection{Appendix B: Benchmark Task Suite (tasks.json, nominal
subset)}\label{appendix-b-benchmark-task-suite-tasks.json-nominal-subset}

\begin{Shaded}
\begin{Highlighting}[]
\OtherTok{[}
  \FunctionTok{\{}
    \DataTypeTok{"id"}\FunctionTok{:} \StringTok{"write\_file"}\FunctionTok{,}
    \DataTypeTok{"prompt"}\FunctionTok{:} \StringTok{"write \textquotesingle{}hello world\textquotesingle{} to /tmp/apex{-}bench{-}out.txt"}\FunctionTok{,}
    \DataTypeTok{"check"}\FunctionTok{:} \StringTok{"hello"}\FunctionTok{,}
    \DataTypeTok{"check\_file"}\FunctionTok{:} \StringTok{"/tmp/apex{-}bench{-}out.txt"}
  \FunctionTok{\}}\OtherTok{,}
  \FunctionTok{\{}
    \DataTypeTok{"id"}\FunctionTok{:} \StringTok{"shell\_echo"}\FunctionTok{,}
    \DataTypeTok{"prompt"}\FunctionTok{:} \StringTok{"run the shell command: echo apex{-}bench{-}ok"}\FunctionTok{,}
    \DataTypeTok{"check"}\FunctionTok{:} \StringTok{"apex{-}bench{-}ok"}
  \FunctionTok{\}}\OtherTok{,}
  \FunctionTok{\{}
    \DataTypeTok{"id"}\FunctionTok{:} \StringTok{"read\_file"}\FunctionTok{,}
    \DataTypeTok{"prompt"}\FunctionTok{:} \StringTok{"write \textquotesingle{}readtest\textquotesingle{} to /tmp/apex{-}bench{-}read.txt then read it back and output the contents"}\FunctionTok{,}
    \DataTypeTok{"check"}\FunctionTok{:} \StringTok{"readtest"}
  \FunctionTok{\}}\OtherTok{,}
  \FunctionTok{\{}
    \DataTypeTok{"id"}\FunctionTok{:} \StringTok{"memory\_write\_read"}\FunctionTok{,}
    \DataTypeTok{"prompt"}\FunctionTok{:} \StringTok{"store the value \textquotesingle{}bench{-}value\textquotesingle{} under key \textquotesingle{}bench{-}key\textquotesingle{} in memory, then retrieve it"}\FunctionTok{,}
    \DataTypeTok{"check"}\FunctionTok{:} \StringTok{"bench{-}value"}
  \FunctionTok{\}}\OtherTok{,}
  \FunctionTok{\{}
    \DataTypeTok{"id"}\FunctionTok{:} \StringTok{"http\_get"}\FunctionTok{,}
    \DataTypeTok{"prompt"}\FunctionTok{:} \StringTok{"make an HTTP GET request to https://httpbin.org/get and output the response"}\FunctionTok{,}
    \DataTypeTok{"check"}\FunctionTok{:} \StringTok{"httpbin"}
  \FunctionTok{\}}
\OtherTok{]}
\end{Highlighting}
\end{Shaded}

The full 12-task suite including adversarial cases is maintained at
\texttt{apex/benchmarks/tasks.json} in the APEX repository.
